\documentclass[14pt]{beamer}
\usetheme{Madrid}
\usecolortheme{default}

\usepackage{graphicx}
\usepackage{float}
\graphicspath{ {.//obrazky/} }

\usepackage[czech]{babel}
\usepackage[T1]{fontenc}
\usepackage[utf8]{inputenc}
\usefonttheme{serif}
\hypersetup{pdfstartview={Fit}}

\title{Zařízení pro snímání gest}
\author{Štěpán Bílek}
\date{2025/2026}

\begin{document}

\begin{frame}
  \titlepage

  \vspace{0.5cm}

  {\footnotesize
    vedoucí práce: Mgr.~Zbyšek Nechanický\\
    oponent: Jaroslav Kořínek\\
  }
\end{frame}

\begin{frame}{Obsah}
  \tableofcontents
\end{frame}

\section{Představení mé osoby}
\begin{frame}{Zájmy}
	\begin{itemize}
    \item programování
    \item technika a konstruování
    \item hra na kytaru
    \item sport - lezení, jízda na kole
  \end{itemize}
\end{frame}

\begin{frame}{Relevantní zkušenosti}
	\begin{block}{Stáž}
    roční studenstká stáž na Fyzikálním ústavu AV ČR
  \end{block}
  \begin{block}{Práce v ČT}
    programování webových aplikací v České televizi
  \end{block}
\end{frame}

\section{Představení práce}
\begin{frame}{Představení práce - úvod}
  \begin{itemize}
    \item ovládat software nebo robota pohyby
    \item nositelné zařízení
    \item spolupráce se spolužákem Lukášem Karáskem
  \end{itemize}
\end{frame}

\begin{frame}{Představení práce - kritéria}
	\begin{itemize}
    \item opravdová využitelnost zařízení
    \item efektní a příjemné používání
    \item rozšiřitelnost pro různé účely
  \end{itemize}
\end{frame}

\begin{frame}{Představení práce - návrh řešení}
	\begin{itemize}
    \item nositelné senzorické zařízení
    \item přenos nezpracovaných dat
    \item ovládací program na počítač
  \end{itemize}
\end{frame}

\begin{frame}{Představení práce - cíle}
  \begin{itemize}
    \item ovládaný program - prezentace
    \item rozpoznávání 2 gest - \uv{doleva} a \uv{doprava}
    \item cíl maturitní práce - posouvání snímků gesty
  \end{itemize}
\end{frame}


\section{Popis vlastního řešení}
\begin{frame}{Popis vlastního řešení - tištěný spoj}
	\begin{itemize}
    \item udělat zařízení co nejjednodušší
    \item využít existující integrované spoje
    \item komponenty připájené napřímo
  \end{itemize}
\end{frame}

\begin{frame}{Popis vlastního řešení - neuronová síť}
	\begin{itemize}
    \item rozpoznání gesta
    \item univerzálnost a rozšiřitelnost
    \item převážně studium teorie
    \item implementace pomocí knihoven
  \end{itemize}
\end{frame}

\begin{frame}{Popis vlastního řešení - ovládací program}
	\begin{itemize}
    \item Python ekosystém
    \item přenos dat přes UDP
    \item trénování modelu UI
    \item uživatelské rozhraní v terminálu
  \end{itemize}
\end{frame}

\begin{frame}{Popis vlastního řešení - testování}
	\begin{itemize}
    \item špatný kontakt
    \item chybné použití knihovny (gesta pouze \uv{doprava})
    \item upravování trénovacích dat (délka, nesprávný model)
  \end{itemize}
\end{frame}


\section{Ukázka vlastní práce}
\begin{frame}{Ukázka vlastní práce - tištěný spoj}
  \begin{columns}
    \column{0.5\textwidth}
    \begin{figure}[H]
      \centering
      \includegraphics[width=\textwidth]{navrh_pcb.png}
    \end{figure}

    \column{0.5\textwidth}
    \begin{itemize}
      \item 2 vrstvy
      \item vlastní obtisky součástek
    \end{itemize}
  \end{columns}
\end{frame}

\begin{frame}{Ukázka vlastní práce - tištěný spoj}
	\begin{columns}
    \column{0.5\textwidth}
    \begin{itemize}
      \item hlavní vypínač
      \item 3+1 tlačítka
      \item 2 LED
      \item 2 konektory
    \end{itemize}

    \column{0.5\textwidth}
    \begin{figure}[h]
      \centering
      \includegraphics[width=0.8\textwidth, angle=180]{neo.jpg}
    \end{figure}
  \end{columns}
\end{frame}

\begin{frame}{Ukázka vlastní práce - tištěný spoj}
	\begin{columns}
    \column{0.5\textwidth}
    \begin{figure}[h]
      \centering
      \includegraphics[width=\textwidth, angle=270]{elek_celek.jpg}
    \end{figure}

    \column{0.5\textwidth}
    \begin{itemize}
      \item Li-pol akumulátor
      \item regulátor nabíjení/vybíjení
    \end{itemize}
  \end{columns}
\end{frame}

\begin{frame}{Ukázka vlastní práce - neuronová síť}
	\begin{columns}
    \column{0.5\textwidth}
    \begin{itemize}
      \item 2 skryté vrstvy
      \item klasifikace 2 tříd
      \item závislý počet vstupů
    \end{itemize}

    \column{0.5\textwidth}
    \begin{figure}[h]
      \centering
      \includegraphics[width=\textwidth]{NN_python.png}
    \end{figure}
  \end{columns}
\end{frame}

\begin{frame}{Ukázka vlastní práce - rozhraní}
	\begin{figure}[h]
    \centering
    \includegraphics[width=\textwidth]{UI.png}
  \end{figure}
\end{frame}

\begin{frame}{Ukázka vlastní práce - rozhraní}
	\begin{columns}
    \column{0.5\textwidth}
    \begin{figure}[h]
      \centering
      \includegraphics[width=\textwidth]{UI_live.png}
    \end{figure}

    \column{0.5\textwidth}
    \begin{itemize}
      \item živá data
    \end{itemize}
  \end{columns}
\end{frame}

\begin{frame}{Ukázka vlastní práce - rozhraní}
	\begin{columns}
    \column{0.5\textwidth}
    \begin{itemize}
      \item stavové a chybové hlášky
      \item soubory pro tréninková data
    \end{itemize}

    \column{0.5\textwidth}
    \begin{figure}[h]
      \centering
      \includegraphics[width=\textwidth]{UI_msg.png}
    \end{figure}
  \end{columns}
\end{frame}

\begin{frame}{Ukázka vlastní práce - rozhraní}
	\begin{figure}[h]
    \centering
    \includegraphics[width=\textwidth]{UI_comm.png}
  \end{figure}
  \begin{itemize}
    \item funkce programu
    \item klávesové zkratky
  \end{itemize}
\end{frame}

\begin{frame}{Ukázka vlastní práce - firmware}
	\begin{itemize}
    \item vyčítání dat z gyroskopu a akcelerometru
    \item vlastní WiFi síť
    \item odesílaní dat po zmáčknutí tlačítka
    \item LED signalizující vysílání
  \end{itemize}
\end{frame}

\section{Závěr}
\begin{frame}{Závěr}
  \begin{block}{Naplnění cílů}
    \begin{itemize}
      \item koncept ověřen
      \item funkční prototyp
    \end{itemize}
  \end{block}

  \begin{block}{Spolupráce}
    \begin{itemize}
      \item čekání na postupu spolupracovníka
      \item obecně vpořádku
    \end{itemize}
  \end{block}
\end{frame}

\begin{frame}{Závěr}
	\begin{block}{Pokračování v práci}
    \begin{itemize}
      \item tištěný spoj M2
      \item automatické rozpoznání gest
    \end{itemize}
  \end{block}
\end{frame}

\section{Časová náročnost}
\begin{frame}{Graf časové náročnosti}
  \begin{figure}[h]
    \centering
    \includegraphics[width=0.6\textwidth]{hodiny.png}
  \end{figure}
  \begin{itemize}
    \item celkový počet hodin: 92
  \end{itemize}
\end{frame}

\end{document}
